\subsection{Irreducible tame residual image leaves only joint inertia eighteen}

\begin{proposition}[B-RD38: the cyclic-four residual case]
\label{prop:B-RD38}
\status{Proved. Complete root and independent proof, original-source, and exact-certificate audits passed.}
Let $K$ be a complete field of mixed characteristic $(0,3)$ with
algebraically closed residue field and absolute ramification index two.
Suppose an elliptic curve $E/K$ has residual image
$\rho_{E,3}(G_K)=C_4\subset\operatorname{SL}_2(\mathbf F_3)$.
Then
\[
                  v_K(j(E))=6,
       \qquad     6\mid v_K(\Delta(E)),
                                                        \tag{B-RD38.1}
\]
where the discriminant is that of any Weierstrass equation.
In particular $E$ is potentially supersingular and its tame
special-inertia quotient has order one or two, never four.

Consequently, in the remaining level-$27$ test of
Theorem~\ref{thm:B-RD34}, a common irreducible $C_4$ residual image
can occur with different prime-to-three characters only if the two
individual faithful special-inertia orders are $3$ and $6$, and the
joint faithful order is $18$. Every possible failure with joint order
$36$ has nontrivial unipotent residual image $C_3$ after the common
quadratic twist in Proposition~\ref{prop:B-RD37}.
\end{proposition}
\begin{proof}
Write $L=K(E[3])$. Since its Galois group has order four, $L/K$
is totally tamely ramified and cyclic. Let $F$ be its quadratic
subfield and let $\sigma$ be the nontrivial element of
$\operatorname{Gal}(F/K)$. The residual image of $G_F$ is the
scalar subgroup $\{\pm I\}$, with nontrivial character given by
$L/F$.

We first identify the projective level-three parameter carefully.
The full-level modular curve $X(3)$ carries the action of
$\operatorname{SL}_2(\mathbf F_3)/\{\pm I\}=A_4$, with quotient
the $j$-line, as in \citet[Section~2, p.~4]{Fisher2015Nine}.
Use the parameter and universal elliptic family in
\citet[Lemma~4.2(i), p.~8]{Fisher2015Nine}:
\[
 E_t:\quad y^2+3txy+(t^3-1)y=x^3,
 \qquad
 j(t)=\frac{27t^3(t^3+8)^3}{(t^3-1)^3}.
                                                        \tag{B-RD38.2}
\]
The nontrivial projective residual image is a double transposition
in $A_4$. After changing the symplectic basis, it acts by
\[
                  h(t)=\frac{t+2}{t-1}.
                                                        \tag{B-RD38.3}
\]
Here no transformation formula is being assumed without a check:
$h^2=1$, $j\circ h=j$, and $h$ interchanges the cusp pairs
$\{\infty,1\}$ and $\{\zeta,\zeta^2\}$, where $\zeta\in K$
is a primitive cube root of unity. Thus $h$ is one of the three
double transpositions, which are conjugate under $A_4$.

A symplectic basis of $E[3]$ over $L$ gives a noncuspidal point
of $X(3)$. The central scalar $-I$ acts trivially on this modular
curve, so its coordinate $t$ belongs to $F$, and
\[
                         \sigma(t)=h(t).
                                                        \tag{B-RD38.4}
\]
This permits $t\in K$ when the point is fixed by $h$; the ensuing
norm identity still holds. Put $z=t-1$. Equation (B-RD38.4)
becomes $\sigma(z)=3/z$, whence $N_{F/K}(z)=3$.

Choose $\pi\in K$ with $\pi^2=3$, as in
Proposition~\ref{prop:B-RD36}. With normalized valuations,
$v_F(3)=4$ and $v_F(\pi)=2$. Since $F/K$ is totally ramified,
$v_K(N_{F/K}(z))=v_F(z)$, so $v_F(z)=2$.
The unit $u=z/\pi$ has norm one. The residue action of
$\sigma$ is trivial, and therefore $\overline u^2=1$:
\[
                         \overline u=1\quad\hbox{or}\quad-1.
                                                        \tag{B-RD38.5}
\]
On the other hand, $(\zeta-1)^2=-3\zeta$ gives
\[
             \left(\frac{\zeta-1}{\pi}\right)^2=-\zeta.
\]
Its residual square is $-1$, so its residue is neither $1$
nor $-1$. The same holds for $(\zeta^2-1)/\pi$.
It follows that the three quantities $t-1,t-\zeta,t-\zeta^2$
all have $F$-valuation two. Thus
\[
 v_F(t^3-1)=6,\qquad v_F(t^3+8)=6,
 \qquad v_F(\Delta(E_t))=30,
 \qquad v_F(j(t))=12.                                  \tag{B-RD38.6}
\]
For the second equality use $t^3+8=(t^3-1)+9$ and $v_F(9)=8$.
Since $j(t)=j(E)\in K$, the last equality gives $v_K(j(E))=6$.

It remains to compare the discriminants of the actual elliptic
curves, rather than just their common modular point. Over $F$,
twist $E_F$ by the quadratic character of $L/F$. This kills its
scalar residual character and produces a full rational level-three
basis. Its point in $X(3)$ is exactly the point $t$ above: the
quadratic twist multiplies the basis descent by a central sign,
which is invisible in $X(3)$. The full-level moduli description
therefore identifies this twist with $E_t$ over $F$.

The quadratic extension $L/F$ is tame and ramified. It has the
form $F(\sqrt d)$ with $v_F(d)$ odd; every unit is a square by
Hensel lifting over the algebraically closed residue field.
A quadratic twist changes the discriminant by $d^6$, up to a
twelfth power from a Weierstrass coordinate change. Hence, writing
$\delta=v_K(\Delta(E))$, equation (B-RD38.6) gives
\[
                   30\equiv2\delta+6\pmod {12},
       \qquad\text{and thus}\qquad 6\mid\delta.
                                                        \tag{B-RD38.7}
\]
This comparison also covers the special values of $j$: the
full-level structure eliminates elliptic automorphisms fixing the
basis, so the identified twist is the same marked object, not
merely a curve with the same $j$-invariant.

The strict positivity of $v_K(j)$ gives potential supersingularity
by Proposition~\ref{prop:B-RD36}, whose discriminant formula now
makes the tame special-inertia order one or two. In a remaining
level-$27$ failure, Theorem~\ref{thm:B-RD34} permits only joint
orders eighteen or thirty-six. Joint order thirty-six would require
tame order four on one side, so it is excluded here. Joint order
eighteen with distinct tame characters forces tame orders one and
two and independent wild order-three images. The individual
faithful orders are therefore three and six. Finally apply the two
residual alternatives in Proposition~\ref{prop:B-RD37}.
\end{proof}

\paragraph{Exact finite certificate.}
The retained program \path{code/verify_c4_residual_gate.py} checks
the modular involution, its cusp permutation, the exact $c_4$ and
discriminant transformation factors, the two normalized residues in
$\mathbf F_9$, the quadratic-twist discriminant identity, and the
valuation congruence (B-RD38.7). The author run passed in
$0.07729700100026093$ seconds; its complete output is retained
beside the program.
These checks establish the finite identities used above; neither
existence nor nonexistence of a remaining $C_3$ or $C_4$ level-$27$
comparison is inferred. Independent author-output-preserving reruns by
the root and second reviewers passed in $0.1112131780027994$ and
$0.10795447899727151$ seconds, respectively.

\paragraph{Subsequent exact determination.}
The remaining small-height alternatives isolated here are settled in
Theorem~\ref{thm:B-RD42}: the exact manuscript index is
$\nu_{\min}(3,2,2)=2$. The intermediate bounds and conditional
reductions above are retained as steps in that proof.
